% INTEGRATED_REVISION_V2
\documentclass[11pt,a4paper]{article}
\usepackage[margin=25mm]{geometry}
\usepackage[T1]{fontenc}
\usepackage{mathpazo,amsmath,amssymb,graphicx,booktabs,longtable,array,caption,setspace,microtype}
\usepackage{placeins}
\usepackage[hidelinks]{hyperref}
\usepackage{fancyhdr}
\usepackage[mathlines]{lineno}
\setstretch{1.2}
\setlength{\parskip}{4pt}
\setlength{\parindent}{0pt}
\captionsetup{font=small,labelfont=bf}
\pagestyle{fancy}\fancyhf{}\fancyhead[L]{\small GBD 2023: infertility contrasts}\fancyhead[R]{\small Additional file 1}\fancyfoot[C]{\thepage}
\setlength{\headheight}{14pt}
\setlength{\emergencystretch}{2em}
\newcommand{\ASPR}{\mathrm{ASPR}}
\title{Additional file 1: Supplementary methods and results}
\author{Companion to the Reproductive Health manuscript draft}
\date{25 October 2026}
\begin{document}\maketitle
\subsection*{Additional arithmetic definitions}
The same-population benchmark is $B_T=\sum_a P_{C,aT}(r_{C,aT}-r_{R,aT})/10^5$. It uses China's actual age distribution rather than standard weights and is not preventable burden. Age-specific contributions to the standardized difference are $w_a(r_{C,at}-r_{R,at})$. The seven published unnormalized standard weights, in ascending age order, are 0.082891, 0.078012, 0.075914, 0.073217, 0.068281, 0.061474 and 0.055113; they sum to 0.494902 before renormalization. References and outcome definitions are given in the main manuscript.
\section{Supplementary methods and results}
\renewcommand{\thetable}{S\arabic{table}}\setcounter{table}{0}
\renewcommand{\thefigure}{S\arabic{figure}}\setcounter{figure}{0}
\subsection{Decomposition of case-count changes}
We expressed the case count as $F(N,\boldsymbol f,\boldsymbol p)=N\sum_a f_a p_a$, where $N$ is the total female population aged 15--49, $f_a$ is its age distribution, and $p_a$ is age-specific prevalence as a proportion. For each factor, we averaged its marginal contribution across all six possible replacement orders between the starting and ending states. These Shapley contributions sum exactly to the observed modeled case-count change. We examined 1990--2023 and 2010--2023 for both causes in all 19 countries. Population size here refers specifically to females aged 15--49; age composition refers to redistribution within this range. The age-specific prevalence component includes epidemiological change, ascertainment and modeling changes and cannot be interpreted as a causal risk-factor contribution.

\subsection{Trends, sensitivities and forecast validation}
We fitted continuous piecewise linear models to log standardized rates, searching integer-year knot locations with zero to two knots and at least five observations per segment, counting shared boundary years in each segment. Model selection used the Bayesian information criterion, including a penalty for knot locations. Segment-specific annual percentage change was $100[\exp(\beta)-1]$; average annual percentage change (AAPC) was obtained from duration-weighted log slopes. These fits summarize modeled trajectories and are not outputs of the proprietary Joinpoint program. AR(1) generalized least-squares refits retained the OLS-selected knots; allowing up to three knots and restricting the period to 2010--2023 were additional sensitivities. Trend confidence intervals are conditional on selected knots and fixed GBD point estimates.

Reference sensitivities used age-specific country population shares fixed at 1990, excluded one reference country at a time, and restricted ages to 20--49 with renormalized weights. Following inspection of the terminal-year endometriosis increase, we added analyses ending in 2021 and 2022; these endpoint analyses were not prespecified. To evaluate numerical sensitivity to published weight rounding, we evaluated all $2^7$ combinations of perturbing unnormalized weights by $\pm0.5\times10^{-6}$. The resulting deterministic envelopes are not uncertainty intervals.

An exploratory forecasting module evaluated each Chinese age-specific prevalence series on the logit scale. Candidates were a constant-last-rate baseline, ARIMA with $p,q\in\{0,1,2\}$ and $d\in\{0,1\}$ selected by training-set corrected AIC, and additive damped-trend exponential smoothing. ARIMA candidates failing convergence were excluded. Rolling training windows contained at least 20 years and assessed horizons 1, 3 and 5; a separate last-10-year holdout assessed horizons 1--10. We specified an exploratory retention rule before running this module: at least 10\% lower mean absolute error (MAE) than the baseline in every age group in both validation schemes, with complete predictions. Future counts would additionally require compatible population projections; none were generated when validation failed. This module was not a validation against independent observed infertility data.

\begin{table}[htbp]\centering\small
\caption{Shapley decomposition of changes in China's modeled prevalent cases, thousands.}\label{tab:decomp}
\setlength{\tabcolsep}{4pt}
\begin{tabular}{llrrrr}\toprule
Cause & Period & Net change & Population size & Age composition & Prevalence\\\midrule
\input{tables/decomposition.tex}\end{tabular}
\caption*{Population size is the female population aged 15--49; age composition is its within-range distribution; prevalence is age-specific modeled infertility prevalence. Contributions may exceed the net change and need not have the same sign. Rounding can cause small displayed discrepancies.}
\end{table}

\begin{table}[htbp]\centering\small
\caption{Restricted-age standardized YLD rates and YLD totals.}\label{tab:yld}
\setlength{\tabcolsep}{3.5pt}
\begin{tabular}{llrrrrrr}\toprule
 & & \multicolumn{2}{c}{YLDs, thousands} & \multicolumn{2}{c}{Rate per 100,000} & & \\
Cause & Year & China & Reference & China & Reference & $D$ & $RR$\\\midrule
\input{tables/yld.tex}\end{tabular}
\caption*{All values are derived point estimates. YLD: years lived with disability.}
\end{table}
\begin{table}[htbp]\centering\small
\caption{Endpoint sensitivity added after inspection of terminal-year estimates.}\label{tab:endpoint}
\begin{tabular}{llrrrrr}\toprule
Cause & End year & China ASPR & Ref. ASPR & $D$ & $RR$ & $D_T-D_{1990}$\\\midrule
\input{tables/endpoint.tex}\end{tabular}
\caption*{The source release remains GBD 2023 in all rows. Ending in 2021 is not substitution of GBD 2021 estimates. Rate units: per 100,000 women.}
\end{table}
\begin{table}[htbp]\centering\small
\caption{Conditional trend sensitivity for prevalence, 1990--2023.}\label{tab:trends}
\begin{tabular}{lllrr}\toprule
Location & Cause & Model & AAPC (\%) & Conditional 95\% CI\\\midrule
\input{tables/trends.tex}\end{tabular}
\caption*{Intervals condition on selected knots and fixed GBD point estimates. AR(1) retains the knots selected by OLS with up to two knots. These intervals exclude GBD, knot-selection and structural model uncertainty. Full periods, measures, knots and residual-correlation estimates are in the accompanying CSV.}
\end{table}
\begin{table}[htbp]\centering\small
\caption{Age-specific forecast-validation MAE ratios relative to the naive baseline.}\label{tab:forecast}
\begin{tabular}{llccc}\toprule
Cause & Model & Rolling: min--max & Last 10 years: min--max & Retained\\\midrule
\input{tables/forecast.tex}\end{tabular}
\caption*{Ranges span seven age groups. Values below 1 favor the candidate. Retention required every age-specific ratio to be below 0.9 in both schemes, with complete predictions. The rule was specified before running this exploratory module, not externally preregistered. Rolling MAE combines horizons 1, 3 and 5; the holdout combines horizons 1--10.}
\end{table}
\subsection{Acquisition and computation audit}
The official query identifiers were:
\begin{description}
\item[Burden:] \texttt{1b13e3376135ab41a4bc5e0ffc60e84f}.
\item[Population:] \texttt{916826fb285f8d3c3e53117dd2133406}.
\item[Pilot:] \texttt{0daed0e53f5ad8ebdf0dc00c23ccd110}.
\item[Annual strata:] \texttt{416daf735b9fc79bec35506514424e73}.
\item[Endpoint strata:] \texttt{76c5523c52f32caf342843daf2326fa4}.
\end{description}
Raw uncertainty limits remain in the original CSVs. Country and outcome mappings, source hashes, rounding sensitivities and all analytical tables are separate machine-readable files. Nine automated core checks plus the new reference-composition accounting checks cover core accounting identities, the complete data grid, denominator consistency and outcome gates; temporary artificial test fixtures were not used as research data.

The protocol was written before analysis but was not externally preregistered. The endpoint and weight-rounding sensitivities were added following quality inspection. Forecasting remained exploratory and was withheld after failed comparative validation. Primary/secondary stratification was completed using annual Number exports and independent endpoint Number/Rate checks. The symmetric reference-composition and age-by-type analyses were added during manuscript review and were not prespecified. The separately acquired NSFG respondent analysis was completed post hoc and independently numerically checked; these survey estimates do not validate GBD cause attribution. Future population scenarios were not performed. The annual stratum rate reconstruction reduced redundant metrics within the remaining GBD download allowance.

\subsection{GATHER reporting map}
This map applies GATHER to the secondary analysis; upstream input reconstruction and primary disease estimation belong to the original GBD study. It is not a claim of independent GBD model replication.
\small
\begin{longtable}{>{\raggedright\arraybackslash}p{0.08\textwidth}>{\raggedright\arraybackslash}p{0.27\textwidth}>{\raggedright\arraybackslash}p{0.54\textwidth}}
\toprule Item & Topic & Location or applicability\\\midrule\endhead
1 & Indicators and populations & Methods: causes, ages, female sex, fixed countries and 1990--2023.\\
2 & Funding & Declarations: author-supplied information remains pending.\\
3 & Data identification and access & Official Results Tool, download date, query identifiers and cited GBD methods.\\
4 & Eligibility and exclusions & Exact outcome filters and completeness gates; no country removals from the main set. Sensitivity exclusions are explicit.\\
5 & Source characteristics & GBD model outputs, dimensional mappings and acquisition audit. Original source-study characteristics are documented by GBD.\\
6 & Potential input biases & Discussion: sparse evidence, diagnostic heterogeneity, attribution and modeled denominators.\\
7 & Other inputs & Matching GBD population and published standard weights with provenance.\\
8 & Input data access & CSVs retained in the research archive; source institution IHME and exact queries identified for reproducibility.\\
9 & Conceptual overview & Equations for within-age pooling, standardization, contrasts and decomposition.\\
10 & Analytical detail & Normalization, validation, comparison, decomposition, trend and sensitivity methods and scripts.\\
11 & Candidate evaluation and selection & Segmented-model BIC, training AICc and explicit forecast-retention criterion.\\
12 & Evaluation results & Reference, endpoint, AR(1), rounding and forecast-validation results.\\
13 & Uncertainty methods and scope & Joint draws unavailable; no propagated derived intervals; conditional trend interval scope stated.\\
14 & Source code access & GitHub repository, pinned dependencies and executable scripts.\\
15 & Machine-readable estimates & Complete derived CSV tables accompany the paper.\\
16 & Quantitative uncertainty & Incomplete for derived estimates: marginal source limits retained but cannot substitute for joint uncertainty. Trend intervals are conditional only.\\
17 & Interpretation and prior evidence & Discussion compares prior GBD studies. Release-to-release differences were not estimated.\\
18 & Limitations & Discussion: endpoint, reference, denominator, source, rounding and inference limitations.\\\bottomrule
\end{longtable}
\normalsize
\begin{figure}[p]\centering\includegraphics[width=\textwidth]{figures/figS1_reference_sensitivity.pdf}
\caption{Sensitivity of the 2023 prevalence difference to reference-country exclusions, fixed 1990 country shares within ages, and restriction to ages 20--49. Dashed lines show the main 15--49 comparison. PCOS signs are stable; endometriosis signs vary. All are point estimates.}\label{fig:sensitivity}\end{figure}
\begin{figure}[p]\centering\includegraphics[width=\textwidth]{figures/figS2_yld_trends.pdf}
\caption{Restricted-age standardized YLD rates for both infertility causes, 1990--2023. These are infertility-impairment YLD rates rather than the total disability burden of either underlying disease.}\label{fig:yld}\end{figure}

\begin{figure}[htbp]\centering\includegraphics[width=\textwidth]{figures/fig1_country_comparison_2023.pdf}
\caption{Country-specific prevalence rates in 2023, restricted to ages 15--49. Dashed lines are the pooled reference. Descriptive point estimates only.}\end{figure}
\begin{figure}[htbp]\centering\includegraphics[width=\textwidth]{figures/fig3_age_disparities.pdf}
\caption{Age contrasts and standard-weighted contributions to changes in China minus reference. Contributions are not case counts or causal effects.}\end{figure}
\begin{figure}[htbp]\centering\includegraphics[width=\textwidth]{figures/fig4_decomposition_and_benchmark.pdf}
\caption{China's case-count decomposition and same-population benchmarks. The latter are signed arithmetic contrasts using China's actual age distribution, not preventable cases.}\end{figure}
\FloatBarrier

\subsection{Exploratory NSFG age-standardization and co-diagnosis sensitivity}
This post hoc analysis used the original CDC female respondent PUF ZIP and expanded CSV; both SHA-256 checks matched the archived research-source manifest (full digests in the machine-readable audit). There were 5,586 original respondents; 4,855 were aged 20--49 at interview. The one additional respondent aged 50 at interview but aged 49 at screening is included in the official 4,856-person age-domain tabulation; its inclusion did not materially alter the earlier CDC replication results. Diagnosis responses of 8, 9 and missing were excluded rather than treated as no diagnosis. Surgical sterility (FECUND=1--2) was not counted as impaired fecundity.

For diagnosis group $g$ and age group $a$, let $\widehat{p}_{ga}=\widehat{Y}_{ga}/\widehat{N}_{ga}$ be the survey-weighted broad current fertility problem prevalence, with $\widehat{Y}_{ga}$ the weighted outcome numerator. Let $q_a=\widehat{N}_a/\sum_b\widehat{N}_b$ be the estimated survey-weighted population age distribution among all 20--49-year-olds. The standardized rate is $\widehat R_g=\sum_a q_a\widehat p_{ga}$. First-order variance accounted for variation in both $\widehat p_{ga}$ and the common estimated $q_a$; centered PSU totals were combined within all 20 strata (four clusters each) for survey-design confidence intervals using 60 degrees of freedom. The log prevalence-ratio contrast used the joint linearized covariance. Detailed reproducing code and derived data (without individual respondents) accompany this supplemental analysis.

\begin{table}[htbp]\centering\small
\caption{NSFG 2022--2023: broad fertility-problem differences under age standardization and co-diagnosis exclusions.}\label{tab:nsfg_age_std}
\setlength{\tabcolsep}{5pt}
\begin{tabular}{llrr}\toprule
Diagnosis & Definition & Difference, pp (95\% CI) & Ratio (95\% CI)\\\midrule
\input{tables/nsfg_ageadjusted.tex}
\end{tabular}
\caption*{Survey design-based, unadjusted or directly age-standardized within US women 20--49; all exploratory post hoc comparisons with no multiplicity correction. ``No other diagnosis'' removes respondents reporting the alternate condition. Intervals do not represent GBD posterior uncertainty.}
\end{table}
\begin{figure}[htbp]\centering\includegraphics[width=\textwidth]{figures/nsfg_age_adjusted_forest.pdf}
\caption{Sensitivity of diagnosis and broad fertility-problem prevalence ratios to age standardization and removal of co-diagnosis. The vertical dashed line is the null ratio of 1; 95\% pointwise intervals use the NSFG survey design. This is a US survey association study, not validation of GBD estimates.}\label{fig:nsfg_age_std}\end{figure}

\input{sections/revision_supplement.tex}
\input{sections/endo_endpoint_supplement.tex}
\input{sections/un_population_supplement.tex}
\end{document}
